% title:          Perfect numbers and even exponents
% area:           arithmetic-functions, primes
% methods:        contradiction
% difficulty:
% difficulty-ai:  medium
% status:
% review:         none
% --- history: all optional, leave EMPTY if unknown ---
% origin:         imc
% origin-date:    2014-07-31
% origin-number:  Day 1, Problem 4
% also-in:
% taken-from:
% references:     IMC 2014 (21st IMC, Blagoevgrad), Day 1, Problem 4, proposed by Javier Rodrigo (Universidad Pontificia Comillas) - https://www.imc-math.org.uk/?year=2014&item=problems ; questions https://www.imc-math.org.uk/imc2014/IMC2014-day1-questions.pdf ; official solution https://www.imc-math.org.uk/imc2014/IMC2014-day1-solutions.pdf
% history:        ai
\begin{problem}
Let \(n > 6\) be a perfect number, and let \(n = p_{1}^{e_{1}} \cdots p_{k}^{e_{k}}\) be its prime factorisation with \[1 < p_{1} < \ldots < p_{k}.\] Prove that \(e_{1}\) is an even number.
\\
A number \(n\) is \textit{perfect} if \(s(n) = 2n\), where \(s(n)=\sum_{\mathbb{N}\ni d|_{\mathbb{Z}}n}d\) is the sum of the divisors of \(n\).
\end{problem}
\begin{solution}
Suppose that \(e_{1}\) is odd, contrary to the statement. We know that
\[
s(n) = \prod_{i=1}^{k} \left(\sum_{j=0}^{e_i}p_{i}^{j}\right) = 2n = 2p_{1}^{e_{1}} \cdots p_{k}^{e_{k}}.
\]
Since \(e_{1}\) is an odd number, \(p_{1} + 1\) divides the first factor\footnote{\text{one could also see it as follows. Since:} $$\left\{(1+p_{1})\left(\sum_{j=0}^{\frac{e_{1}-1}{2}}(-p_{1})^j\right),(1-p_{1})\left(\sum_{j=0}^{\frac{e_{1}-1}{2}}(-p_{1})^j\right)\right\}=\left\{\left(1-p_{1}^{\frac{e_1+1}{2}}\right),\left(1+p_{1}^{\frac{e_1+1}{2}}\right)\right\}$$\text{where identifying which one is which one depends on the parity of} $\frac{e_1-1}{2}$\text{. We must have that}
\[
\sum_{j=0}^{e_1}p_{1}^{j}=\frac{1-p_{1}^{e_1+1}}{1-p_{1}}=\frac{\left(1-p_{1}^{\frac{e_1+1}{2}}\right)\left(1+p_{1}^{\frac{e_1+1}{2}}\right)}{1-p_{1}}
=\frac{(1+p_{1})\left(\sum_{j=0}^{\frac{e_{1}-1}{2}}(-p_{1})^j\right)(1-p_{1})\left(\sum_{j=0}^{\frac{e_{1}-1}{2}}(-p_{1})^j\right)}{1-p_{1}}
\]
\[
=(1+p_{1})\left(\sum_{j=0}^{\frac{e_{1}-1}{2}}(-p_{1})^j\right)^2.
\]
\text{Comparing with the previous factorization and using its uniqueness yields:}
$$\left(\sum_{j=0}^{\frac{e_1-1}{2}}p_1^{2j}\right)=\left(\sum_{j=0}^{\frac{e_{1}-1}{2}}(-p_{1})^j\right)^2$$}:
$$\sum_{j=0}^{e_1}p_{1}^{j}=\left(\sum_{j=0}^{\frac{e_1-1}{2}}p_1^{2j}\right)+\left(\sum_{j=0}^{\frac{e_1-1}{2}}p_1^{2j+1}\right)=(1+p_1)\left(\sum_{j=0}^{\frac{e_1-1}{2}}p_1^{2j}\right)$$
so \(p_{1} + 1\) divides \(2n\). Due to \(p_{1} + 1 > 2\), at least one of the primes \(p_{1}, \ldots, p_{k}\) divides \(p_{1} + 1\). The primes \(p_{3}, \ldots, p_{k}\) are greater than \(p_{1} + 1\) and \(p_{1}\) cannot divide \(p_{1} + 1\), clearly $p_{1}$ doesn't divide $p_{1}+1$ (else $p_{1}$ would divide $1$) so \(p_{2}\) must divide \(p_{1} + 1\) i.e. $\exists t\in\mathbb{N}^{*}$ with $p_{1}+1=tp_{2}$. Since \(tp_{2}=p_{1} + 1 < 2p_{1}<2p_{2}\), this is possible only if $t=1$ i.e. \(p_{2} = p_{1} + 1\), therefore \(p_{1} = 2\) and \(p_{2} = 3\) (the only two consecutive primes are $2$ and $3$). Hence, \(6 \mid n\).
\\\\
Now \(n, \frac{n}{2}, \frac{n}{3}, \frac{n}{6}\) and \(1\) are distinct divisors of \(n\), so by definition
\[
2n\overset{n\text{ is perfect}}{=}s(n) \geq n + \frac{n}{2} + \frac{n}{3} + \frac{n}{6} + 1 = n\left(1+\frac{1}{2} + \frac{1}{3} + \frac{1}{6}\right)+1= 2n + 1 > 2n,
\]
which is a contradiction.
\begin{remark}
The perfect numbers with a first power odd must therefore be smaller than $6$. One can check that the only perfect number $n\leq 6$ is $6=2\cdot 3$ and this one has an odd first power.
\end{remark}
\end{solution}
